\section{Few-body continua and CDCC}
\label{sec:few-body-cdcc}
% BEGIN CURATED SUBJECT INDEX
\index{partial-wave expansion}
\index{scattering!Coulomb}
\index{tensor product!gauge constraint}
\index{exact WKB}
\index{radial equation}
\index{complex scaling}
\index{pseudostate}
\index{Feshbach projection}
\index{self-energy}
\index{Jacobi coordinate}
% END CURATED SUBJECT INDEX

\calculationroute
  {Tensor-product channels, Jacobi coordinates, partial waves, projection
  operators, Coulomb asymptotics, and finite-basis diagonalization.}
  {Derive three- and four-body CDCC, construct bin and pseudostate model
  spaces, calculate coupling potentials and breakup amplitudes, and smooth
  discrete spectra with complex scaling.}
  {Write the coupled radial equations from a declared few-body Hamiltonian,
  impose open and closed channel conditions, demonstrate model-space
  convergence, and reconstruct an energy-differential observable.}

A weakly bound projectile turns a two-body collision into a spectral problem
with many nearby thresholds.  The continuum is not a small correction: tidal
and Coulomb fields can polarize or break the projectile while it traverses the
target field.  Continuum-discretized coupled channels (CDCC) is a controlled
way to keep a finite part of that continuum in the reaction space
\cite{Kamimura1986CDCC,Austern1987CDCC,Yahiro2012CDCC}.  The method is often
introduced as a numerical prescription.  Here we derive it as a projection of
the direct-integral channel decomposition from Stage I.

This perspective also connects CDCC to the resonance program.  A Feshbach
operator eliminates the continuum and retains its effect in an energy-
dependent self-energy.  CDCC instead retains a finite-rank approximation to
the important continuum fibers.  Complex scaling supplies both resonance
pseudostates and a smoothing representation of the discarded continuum.
Exact WKB identifies the invariant analytic boundary data that any such finite
representation must approximate.

\subsection{Two-cluster projectile: coordinates and Hamiltonian}
\label{subsec:two-cluster-projectile}
% BEGIN CURATED SUBJECT INDEX
\index{operator!adjoint and self-adjoint}
\index{central decomposition}
\index{analytic dilation}
\index{CDCC}
% END CURATED SUBJECT INDEX

Let a projectile $P=c+v$ consist of a core $c$ and a valence particle $v$,
and let it scatter from a target $T$.  After removing the total center of
mass, choose
\begin{align}
  \bm r&=\bm r_v-\bm r_c,
  \label{eq:cdcc-internal-r}\\
  \bm R&=\frac{m_c\bm r_c+m_v\bm r_v}{m_P}-\bm r_T,
  \qquad m_P=m_c+m_v .
  \label{eq:cdcc-projectile-target-R}
\end{align}
The fragment--target coordinates are
\begin{equation}
  \bm R_c=\bm R-\frac{m_v}{m_P}\bm r,
  \qquad
  \bm R_v=\bm R+\frac{m_c}{m_P}\bm r.
  \label{eq:cdcc-fragment-target-coordinates}
\end{equation}
With intrinsic core and target excitations temporarily frozen, the Hamiltonian
is
\begin{equation}
  H=K_R+h_P+U_{cT}(\bm R_c)+U_{vT}(\bm R_v),
  \label{eq:three-body-reaction-H}
\end{equation}
where
\begin{equation}
  h_P=K_r+V_{cv}(\bm r),
  \qquad
  K_R=-\frac{\hbar^2}{2\mu_R}\nabla_R^2,
  \qquad
  K_r=-\frac{\hbar^2}{2\mu_r}\nabla_r^2.
  \label{eq:projectile-internal-H}
\end{equation}
The optical interactions $U_{xT}$ may be complex because target excitations
not included explicitly have already been eliminated.  This is an effective
non-self-adjointness distinct from analytic dilation.  The full microscopic
Hamiltonian before elimination may still be self-adjoint.

The projectile has a bound ground state $\Phi_0$ and continuum states
$\Phi_{\epsilon\alpha}$,
\begin{align}
  h_P\Phi_0&=\epsilon_0\Phi_0,
  \qquad \epsilon_0<0,\\
  h_P\Phi_{\epsilon\alpha}&=\epsilon\Phi_{\epsilon\alpha},
  \qquad \epsilon\geq0,
  \label{eq:projectile-spectrum}
\end{align}
where $\alpha$ collects orbital angular momentum, spin coupling, and any
rearrangement label.  The exact projectile resolution is
\begin{equation}
  \identity_P
  =\ket{\Phi_0}\bra{\Phi_0}
   +\sum_\alpha\int_0^\infty
    \ket{\Phi_{\epsilon\alpha}}
    \bra{\Phi_{\epsilon\alpha}}\dd\epsilon
   +\sum_{b\neq0}\ket{\Phi_b}\bra{\Phi_b},
  \label{eq:projectile-completeness}
\end{equation}
with the normalization and density convention absorbed into
$\Phi_{\epsilon\alpha}$.  CDCC replaces a physically selected part of the
integral by a finite orthonormal set.

\subsection{Model-space projection and the CDCC equation}
\label{subsec:cdcc-projection}
% BEGIN CURATED SUBJECT INDEX
\index{CDCC}
\index{nuclear breakup}
% END CURATED SUBJECT INDEX

Choose square-integrable projectile states $\{\Phi_i\}_{i=0}^{N}$ and define
\begin{equation}
  P_N=\sum_{i=0}^{N}\ket{\Phi_i}\bra{\Phi_i},
  \qquad Q_N=\identity-P_N.
  \label{eq:cdcc-model-projector}
\end{equation}
The CDCC approximation is not simply $\Psi\approx P_N\Psi$.  It is the
projected boundary-value problem
\begin{equation}
  P_N(H-\vsub{E}{tot})P_N\Psi_N=0
  \label{eq:cdcc-projected-equation}
\end{equation}
with scattering conditions imposed in every retained open channel.  Expand
for total angular momentum $JM$,
\begin{equation}
  \Psi^{JM}_N(\bm R,\xi)
  =\frac{1}{R}\sum_c
  \chi_c^J(R)
  \mathcal Y_c^{JM}(\hat{\bm R},\xi),
  \label{eq:cdcc-channel-expansion}
\end{equation}
where $\xi$ denotes all projectile internal coordinates and
\begin{equation}
  \mathcal Y_c^{JM}
  =\left[\Phi_{iI}(\xi)\otimes Y_L(\hat{\bm R})\right]_{JM},
  \qquad c=(i,I,L,\ldots).
  \label{eq:cdcc-channel-function}
\end{equation}
Projecting onto $\mathcal Y_c^{JM}$ gives
\begin{align}
  &\left[
  -\frac{\hbar^2}{2\mu_R}
   \left(\frac{\rmd^2}{\rmd R^2}-\frac{L_c(L_c+1)}{R^2}\right)
  +\epsilon_c-\vsub{E}{tot}\right]\chi_c^J(R)
  \notag\\
  &\hspace{3cm}
  +\sum_{c'}U^J_{cc'}(R)\chi_{c'}^J(R)=0,
  \label{eq:cdcc-coupled-radial}
\end{align}
with
\begin{equation}
  U^J_{cc'}(R)=
  \Bra{\mathcal Y_c^{JM}}
  U_{cT}+U_{vT}\Ket{\mathcal Y_{c'}^{JM}}_{\xi,\hat R}.
  \label{eq:cdcc-coupling-potential}
\end{equation}
Equation~\eqref{eq:cdcc-coupled-radial} is the central computational object.
The off-diagonal matrix elements contain excitation, breakup, de-excitation,
and continuum--continuum coupling on equal footing.

The discarded space generates the exact correction
\begin{equation}
  \Sigma_N(E)=P_NHQ_N
  \frac{1}{E-Q_NHQ_N+\rmi0}
  Q_NHP_N.
  \label{eq:cdcc-discarded-selfenergy}
\end{equation}
CDCC neglects $\Sigma_N$ only after moving the dominant configurations into
$P_N$.  Convergence is therefore a statement that matrix elements of
$\Sigma_N$ relevant to the desired observables have become small, not that
the continuum has ceased to exist.

\subsection{Bin discretization}
\label{subsec:cdcc-bins}
% BEGIN CURATED SUBJECT INDEX
\index{Hilbert space}
\index{wave packet}
\index{partial-wave expansion}
\index{CDCC}
\index{CDCC!bin discretization}
% END CURATED SUBJECT INDEX

For a two-body projectile, a momentum-bin state in partial wave $\alpha$ is
\begin{equation}
  \vsup{\Phi_{n\alpha}}{bin}
  =\frac{1}{\sqrt{N_{n\alpha}}}
  \int_{k_{n-1}}^{k_n}
  f_{n\alpha}(k)\Phi_{k\alpha}\dd k,
  \label{eq:cdcc-bin-state}
\end{equation}
with
\begin{equation}
  N_{n\alpha}=
  \int_{k_{n-1}}^{k_n}|f_{n\alpha}(k)|^2\dd k.
  \label{eq:cdcc-bin-normalization}
\end{equation}
Nonoverlapping bins are orthonormal if the continuum states are delta
normalized.  A constant $f$ represents a simple average; a phase such as
$\rme^{-\rmi\delta_\alpha(k)}$ can reduce rapid scattering-phase variation
within the bin.  Resonant bins may be narrowed around a phase-shift rise, but
the bin width must remain finite until all continuum-normalized quantities
have been combined.

The finite resolution has a direct-integral meaning.  The characteristic
functions of the bins generate a finite-dimensional subspace of
$L^2([0,k_{\max}],\rmd k)$.  Refining the partition produces strong
convergence to the identity on square-integrable wave packets.  It does not
produce norm convergence on the whole continuum Hilbert space.  This is the
same distinction that makes an individual $\ket{k}$ distributional.

\subsection{Pseudostate discretization and the Gaussian expansion method}
\label{subsec:cdcc-pseudostates}
% BEGIN CURATED SUBJECT INDEX
\index{spectrum!point}
\index{pseudostate}
\index{Jacobi coordinate}
\index{CDCC}
\index{Gaussian expansion method}
\index{CDCC!four-body}
% END CURATED SUBJECT INDEX

For a many-body projectile, exact continuum wave functions are expensive and
carry several momenta.  The pseudostate method diagonalizes $h_P$ in a finite
$L^2$ basis,
\begin{equation}
  \sum_{j=1}^{N_b}
  \bra{\varphi_i}h_P\ket{\varphi_j}c_j^{(n)}
  =\epsilon_n
  \sum_{j=1}^{N_b}\braket{\varphi_i|\varphi_j}c_j^{(n)}.
  \label{eq:pseudostate-generalized-eigenproblem}
\end{equation}
Negative eigenvalues approximate bound states; positive eigenvalues provide
quadrature-like samples of the continuum as seen by the chosen basis and
operators.  They are not physical bound levels.  Their density changes with
the basis and is itself useful information only after smoothing.

The Gaussian expansion method (GEM) uses ranges in geometric progression
\cite{Hiyama:2003cu},
\begin{equation}
  \varphi_{nlm}(\bm r)
  =N_{nl}r^l\rme^{-(r/r_n)^2}Y_{lm}(\hat{\bm r}),
  \quad
  r_n=r_1a^{n-1}.
  \label{eq:gem-basis}
\end{equation}
Short and long ranges are then covered with relatively few functions.  For a
three-cluster projectile, products of Gaussians in several Jacobi coordinate
sets efficiently represent both compact correlations and rearrangement tails.
Complex-range Gaussians can improve oscillatory continuum resolution, but
their overlap matrices require careful conditioning.

\subsection{Open, closed, and Coulomb-distorted channels}
\label{subsec:cdcc-boundary-conditions}
% BEGIN CURATED SUBJECT INDEX
\index{boundary condition!incoming and outgoing}
\index{scattering!Coulomb}
\index{CDCC}
% END CURATED SUBJECT INDEX

For a channel with projectile energy $\epsilon_c$, define
\begin{equation}
  E_c=\vsub{E}{tot}-\epsilon_c,
  \qquad
  K_c=\sqrt{2\mu_RE_c}/\hbar.
  \label{eq:cdcc-channel-momentum}
\end{equation}
If $E_c>0$, the channel is open.  With Coulomb interactions, the asymptotic
condition for incidence in $c_0$ is
\begin{equation}
  \chi_c^J(R)\longrightarrow
  \frac{\rmi}{2}\left[
  H_{L_c}^{(-)}(\eta_c,K_cR)\delta_{cc_0}
  -\sqrt{\frac{K_{c_0}}{K_c}}
   S^J_{cc_0}H_{L_c}^{(+)}(\eta_c,K_cR)
  \right].
  \label{eq:cdcc-open-boundary}
\end{equation}
The normalization factor ensures flux normalization.  For a neutral channel,
$H_L^{(\pm)}$ reduce to spherical Hankel combinations.  This is a concrete
instance of the Stage II lesson: a charged channel is compared with a Coulomb
wave, not a plane wave.

If $E_c<0$, write $K_c=\rmi\kappa_c$ with $\kappa_c>0$ and impose the decaying
Whittaker solution,
\begin{equation}
  \chi_c^J(R)\longrightarrow
  C_cW_{-\eta_c,L_c+1/2}(2\kappa_cR).
  \label{eq:cdcc-closed-boundary}
\end{equation}
Closed channels carry no asymptotic outgoing flux, but they alter open-channel
scattering through virtual transitions.  Omitting them simply because they
are not directly detected can be a large error at low collision energy.

Flux conservation supplies a stringent check when all interactions are real:
\begin{equation}
  \sum_{c\,\mathrm{open}}|S^J_{cc_0}|^2=1.
  \label{eq:cdcc-unitarity}
\end{equation}
For optical interactions the deficit equals absorption into eliminated
channels.  Numerically, one must distinguish that physical deficit from loss
caused by insufficient matching radius or an ill-conditioned basis.

\subsection{Multipole expansion of the coupling}
\label{subsec:cdcc-multipoles}
% BEGIN CURATED SUBJECT INDEX
\index{scattering!Coulomb}
\index{CDCC}
\index{nuclear breakup}
% END CURATED SUBJECT INDEX

The fragment--target potentials depend on both $R$ and internal coordinates.
For a central fragment potential,
\begin{equation}
  U_{xT}(|\bm R+\alpha_x\bm r|)
  =\sum_{\lambda\mu}
  \mathcal U_{x\lambda}(R,r)
  Y_{\lambda\mu}(\hat R)
  Y_{\lambda\mu}^*(\hat r),
  \label{eq:cdcc-multipole-expansion}
\end{equation}
where $\alpha_v=m_c/m_P$ and $\alpha_c=-m_v/m_P$.  Angular momentum algebra
reduces Eq.~\eqref{eq:cdcc-coupling-potential} to radial transition densities
and Clebsch--Gordan, $6j$, and $9j$ coefficients.  The long-range Coulomb
multipoles behave as
\begin{equation}
  \mathcal U^{C}_{\lambda}(R,r)
  \sim \frac{r^\lambda}{R^{\lambda+1}},
  \qquad r<R,
  \label{eq:coulomb-multipole-tail}
\end{equation}
so electric dipole or quadrupole breakup can couple to very large $R$.
Convergence in matching radius and projectile excitation energy must therefore
be tested together.

The diagonal monopole contains the projectile--target Coulomb comparison
potential.  Subtracting it before numerical propagation and restoring its
known Coulomb functions in Eq.~\eqref{eq:cdcc-open-boundary} prevents the
long-range phase from being mistaken for a short-range coupling.

\subsection{Three-cluster projectiles and four-body CDCC}
\label{subsec:four-body-cdcc}
% BEGIN CURATED SUBJECT INDEX
\index{radial equation}
\index{pseudostate}
\index{Jacobi coordinate}
\index{antisymmetrization}
\index{CDCC}
\index{CDCC!four-body}
% END CURATED SUBJECT INDEX

Now let $P=a+b+c$.  The projectile plus target is a four-body reaction
system.  A Jacobi set, for example,
\begin{equation}
  \bm x=\bm r_b-\bm r_a,
  \qquad
  \bm y=\bm r_c-
  \frac{m_a\bm r_a+m_b\bm r_b}{m_a+m_b},
  \label{eq:three-cluster-jacobi}
\end{equation}
gives
\begin{equation}
  h_P=K_x+K_y+V_{ab}+V_{bc}+V_{ca}+V_{3b}.
  \label{eq:three-cluster-projectile-H}
\end{equation}
The optional effective three-body interaction $V_{3b}$ represents correlations
not reproduced by pairwise cluster interactions and is calibrated only within
a declared structure model.

There are three Jacobi rearrangement sets.  A basis restricted to one set may
resolve one clusterization well and another poorly.  GEM therefore uses the
union
\begin{equation}
  \Phi_{nI}=
  \sum_{c=1}^{3}\sum_{\alpha_c}
  C_{n\alpha_c}^{(c)}
  \left[
  \varphi_{n_xl_x}^{(c)}(\bm x_c)
  \otimes
  \varphi_{n_yl_y}^{(c)}(\bm y_c)
  \right]_{I},
  \label{eq:gem-rearrangement-expansion}
\end{equation}
followed by antisymmetrization or Pauli projection as required.  Diagonalizing
$h_P$ generates a ground state, physical resonances represented by stabilized
pseudostates, and nonresonant continuum pseudostates.  Inserting these states
into Eq.~\eqref{eq:cdcc-channel-expansion} gives four-body CDCC
\cite{Matsumoto2004FourBodyCDCC,Matsumoto:2010mi}.

The name counts particles in the reaction Hamiltonian, not the number of
coupled radial equations.  A three-body projectile may yield thousands of
channels after angular momentum coupling.  Sparsity, selection rules, and
model-space diagnostics are therefore part of the physics calculation.

\subsection{Folding interactions and transition densities}
\label{subsec:cdcc-folding}
% BEGIN CURATED SUBJECT INDEX
\index{monodromy}
\index{CDCC}
\index{folding interaction}
% END CURATED SUBJECT INDEX

For nucleon scattering from a composite projectile, a semi-microscopic
coupling can be built by folding an effective nucleon--nucleon interaction
with projectile transition densities.  Schematically,
\begin{equation}
  U^J_{cc'}(R)
  =\sum_{\tau=p,n}
   \int\rmd^3s\,
   \rho^{cc'}_\tau(\bm s)
   g_{N\tau}(\bm R-\bm s;\rho,E),
  \label{eq:single-folding-coupling}
\end{equation}
where
\begin{equation}
  \rho^{cc'}_\tau(\bm s)
  =\bra{\Phi_c}
  \sum_{i\in\tau}\delta(\bm s-\bm s_i)
  \ket{\Phi_{c'}}.
  \label{eq:transition-density}
\end{equation}
The $g$ matrix summarizes in-medium multiple scattering.  A phenomenological
renormalization
\begin{equation}
  g\mapsto N_V\re g+\rmi N_W\im g
  \label{eq:jlm-renormalization}
\end{equation}
must be interpreted as model calibration, not as the monodromy renormalization
of Stage V.  The identical word ``renormalization'' refers to different
operations; naming the varied scale and the invariant observable removes the
ambiguity.  The JLM interaction is a standard example
\cite{Jeukenne1977JLM}.

\subsection{From discrete breakup strengths to a continuous spectrum}
\label{subsec:cdcc-smoothing}
% BEGIN CURATED SUBJECT INDEX
\index{spectrum!continuous}
\index{resolvent}
\index{Lippmann--Schwinger equation}
\index{complex scaling}
\index{biorthogonality}
\index{continuum level density}
\index{pseudostate}
\index{CDCC}
\index{CDCC!four-body}
\index{nuclear breakup}
% END CURATED SUBJECT INDEX

The CDCC equations return transition amplitudes $T_n$ to normalized
pseudostates.  An experiment resolves continuum energy and fragment angles,
so the discrete amplitudes must be embedded back into the continuum.  If
$\Phi_{\epsilon\alpha}^{(-)}$ is the exact projectile scattering state, then
\begin{equation}
  T_\alpha(\epsilon)
  \approx\sum_{n=1}^{N}
  \braket{\Phi_{\epsilon\alpha}^{(-)}|\Phi_n}T_n.
  \label{eq:cdcc-exact-smoothing}
\end{equation}
For a two-cluster projectile these overlaps can be calculated directly.  For
a three-cluster continuum, that approach is as hard as the original few-body
scattering problem.

Complex scaling replaces the oscillatory continuum Green function by an
$L^2$ spectral expansion.  Start from the Lippmann--Schwinger state
\begin{equation}
  \ket{\Phi_\epsilon^{(-)}}
  =\ket{\phi_0}
  +\frac{1}{\epsilon-h_P-\rmi0}V\ket{\phi_0}.
  \label{eq:projectile-lippmann-schwinger}
\end{equation}
For dilation-analytic interactions,
\begin{equation}
  \frac{1}{\epsilon-h_P-\rmi0}
  =U^{-1}(\theta)
   \frac{1}{\epsilon-h_P^\theta}
   U(\theta),
  \label{eq:csls-resolvent-rotation}
\end{equation}
when the rotation exposes the required boundary value.  Expanding the scaled
resolvent in biorthogonal eigenvectors,
\begin{equation}
  \frac{1}{\epsilon-h_P^\theta}
  \approx\sum_\nu
  \frac{\ket{\Phi_\nu^\theta}
        \bra{\widetilde\Phi_\nu^\theta}}
       {\epsilon-\epsilon_\nu^\theta},
  \label{eq:csls-pseudostate-resolvent}
\end{equation}
turns the smoothing overlap into finite $L^2$ matrix elements.  This is the
complex-scaled Lippmann--Schwinger (CSLS) construction
\cite{Kikuchi2009CSLS,Kikuchi2010CSLS}.

For an inclusive energy distribution one may avoid constructing amplitudes
channel by channel.  If $\ket{F}=\sum_nT_n\ket{\Phi_n}$ is the projectile
source produced by the reaction, the response is
\begin{align}
  \frac{\rmd\sigma}{\rmd\epsilon}
  &\propto
  -\frac{1}{\pi}\im
  \bra{F}(\epsilon-h_P+\rmi0)^{-1}\ket{F}
  \notag\\
  &\approx-\frac{1}{\pi}\im\sum_\nu
  \frac{
  \bra{F}U^{-1}(\theta)\ket{\Phi_\nu^\theta}
  \bra{\widetilde\Phi_\nu^\theta}U(\theta)\ket{F}}
  {\epsilon-\epsilon_\nu^\theta}.
  \label{eq:complex-scaled-breakup-response}
\end{align}
This is the same resolvent identity used for continuum level density.  The
reaction determines the source $F$; the intrinsic complex-scaled spectrum
determines how that source is distributed among poles and continuum.

\subsection{Selecting a rearrangement channel}
\label{subsec:cdcc-rearrangement-selection}
% BEGIN CURATED SUBJECT INDEX
\index{cross section}
\index{scattering!Coulomb}
\index{scattering!three-body Coulomb}
\index{pseudostate}
\index{CDCC}
\index{nuclear breakup}
% END CURATED SUBJECT INDEX

For three fragments, ``the breakup cross section'' is not a single asymptotic
channel.  One may observe sequential breakup through a two-body resonance,
direct three-body breakup, or a different cluster partition.  Let
$h_P=h_\alpha+K_\alpha+\vsup{V_\alpha}{rem}$ be adapted to a chosen partition.
The CSLS state constructed with the corresponding channel Hamiltonian
$h_\alpha+K_\alpha$ projects the common CDCC source onto that asymptotic
arrangement.  Symbolically,
\begin{equation}
  T_\alpha(\epsilon,\Omega)
  =\sum_n
  \braket{\Phi_{\epsilon\Omega,\alpha}^{(-)}|\Phi_n}T_n.
  \label{eq:arrangement-resolved-amplitude}
\end{equation}
Summing mutually orthogonal arrangements is safe only when their asymptotic
spaces and overlap regions have been defined so that double counting is
controlled.

This point is where nuclear few-body scattering meets the long-range problem
of Stage II.  With three charged fragments, different cluster regions have
different logarithmic Coulomb phases.  A finite internal pseudostate basis can
calculate observables over a controlled interaction region, but it does not
constitute a proof of one uniform three-body Coulomb asymptotic formula.  The
comparison dynamics and channel atlas must be declared separately.

\subsection{Observables from the channel S matrix}
\label{subsec:cdcc-observables}
% BEGIN CURATED SUBJECT INDEX
\index{cross section}
\index{scattering matrix}
\index{scattering!Coulomb}
\index{pseudostate}
\index{optical potential}
\index{CDCC}
\index{nuclear breakup}
% END CURATED SUBJECT INDEX

For spinless notation, the scattering amplitude from $c_0$ to an open channel
$c$ has the partial-wave form
\begin{equation}
  f_{cc_0}(\Omega)
  =f_C(\Omega)\delta_{cc_0}
  +\frac{2\pi}{\rmi\sqrt{K_cK_{c_0}}}
   \sum_{JMLM_L}
   \mathcal A_{c}^{JM}(\Omega)
   [S^J_{cc_0}-\delta_{cc_0}],
  \label{eq:cdcc-scattering-amplitude}
\end{equation}
where $\mathcal A$ contains the angular momentum coefficients and Coulomb
phases.  The differential cross section contains the flux ratio,
\begin{equation}
  \frac{\rmd\sigma_{c\leftarrow c_0}}{\rmd\Omega}
  =\frac{K_c}{K_{c_0}}|f_{cc_0}(\Omega)|^2.
  \label{eq:cdcc-differential-cross-section}
\end{equation}
For pseudostates, summing Eq.~\eqref{eq:cdcc-differential-cross-section} over a
finite energy interval is meaningful after convergence.  A density per unit
energy requires Eq.~\eqref{eq:cdcc-exact-smoothing} or
Eq.~\eqref{eq:complex-scaled-breakup-response}; dividing a discrete value by
an arbitrary pseudostate spacing is not a basis-independent observable.

The total reaction cross section for entrance wave number $K_0$ is
\begin{equation}
  \sigma_R=\frac{\pi}{K_0^2}
  \sum_J(2J+1)
  \left[1-\sum_{c\,\mathrm{open}}|S^J_{cc_0}|^2\right]
  \label{eq:cdcc-reaction-cross-section}
\end{equation}
in a simple spin convention.  With real couplings the bracket vanishes after
all retained open channels are included.  With absorptive optical potentials
it measures loss into channels outside the explicit CDCC space.

\subsection{CDCC as a spectral quadrature}
\label{subsec:cdcc-spectral-quadrature}
% BEGIN CURATED SUBJECT INDEX
\index{spectral measure}
\index{spectrum!point}
\index{complex scaling}
\index{pseudostate}
\index{CDCC}
% END CURATED SUBJECT INDEX

The direct-integral view gives a compact mathematical statement of the
method.  Let $P(\Delta)$ be the projectile spectral measure.  For functions
of $h_P$ tested between localized form factors,
\begin{equation}
  \bra{f}g(h_P)P([0,E_{\max}])\ket{g}
  =\int_0^{E_{\max}}g(\epsilon)\dd\mu_{fg}(\epsilon).
  \label{eq:cdcc-spectral-integral}
\end{equation}
A pseudostate model replaces this by a quadrature
\begin{equation}
  \int_0^{E_{\max}}g(\epsilon)\dd\mu_{fg}(\epsilon)
  \approx\sum_{n=1}^{N}
  g(\epsilon_n)
  \braket{f|\Phi_n}\braket{\Phi_n|g}.
  \label{eq:cdcc-spectral-quadrature}
\end{equation}
Accuracy is therefore operator and source dependent.  A model space that
converges an elastic amplitude need not resolve a sharply differential
three-fragment observable.  Conversely, pseudostate eigenvalues need not
approximate every point of the continuum uniformly for the tested integral to
converge rapidly.

Complex scaling deforms the quadrature contour.  Stable discrete poles remain
fixed while the nonresonant quadrature points follow the rotated cut.  This is
precisely the pole-versus-representation separation established in Stage IV.

\subsection{Connection data and coupled-channel exact WKB}
\label{subsec:cdcc-exact-wkb-connection}
% BEGIN CURATED SUBJECT INDEX
\index{spectrum!point}
\index{scattering matrix}
\index{scattering!multichannel}
\index{boundary condition!incoming and outgoing}
\index{exact WKB}
\index{adiabatic approximation}
\index{CDCC}
\index{holonomy}
% END CURATED SUBJECT INDEX

For $N$ retained channels, Eq.~\eqref{eq:cdcc-coupled-radial} is a matrix
Schr\"odinger equation.  Fundamental solution matrices replace scalar WKB
branches, and the asymptotic coefficients form vectors.  A global transfer
matrix has blocks
\begin{equation}
  \begin{pmatrix}
  \vsub{\bm a}{out}^{R}\\
  \vsub{\bm a}{in}^{R}
  \end{pmatrix}
  =\mathcal M(E)
  \begin{pmatrix}
  \vsub{\bm a}{out}^{L}\\
  \vsub{\bm a}{in}^{L}
  \end{pmatrix}.
  \label{eq:multichannel-transfer}
\end{equation}
A resonance satisfies the determinant condition
\begin{equation}
  \det\Cincoming(E)=0,
  \label{eq:multichannel-resonance-determinant}
\end{equation}
where $\Cincoming$ is the block multiplying prohibited incoming
components.  CDCC approximates this matrix connection problem on the real
axis and can be analytically continued or complex scaled.  Exact WKB supplies
a nonperturbative construction when the matrix spectral curve and its Stokes
data can be controlled.  The scalar theory of Stage III is therefore the
first member of a coupled-channel hierarchy rather than a separate subject.

The eigenvalues of the local potential matrix are sometimes called adiabatic
surfaces.  Their complex degeneracies generate nonadiabatic connection
matrices.  At a coalescence of resonant solutions,
Eq.~\eqref{eq:multichannel-resonance-determinant} has a multiple zero and the
exceptional-point analysis of Stage VII applies.  Few-body reaction theory,
exact WKB, and non-Hermitian holonomy meet at the same determinant.

\subsection{A reproducible convergence protocol}
\label{subsec:cdcc-convergence-protocol}
% BEGIN CURATED SUBJECT INDEX
\index{partial-wave expansion}
\index{scattering!Coulomb}
\index{pseudostate}
\index{Jacobi coordinate}
\index{CDCC}
\index{folding interaction}
% END CURATED SUBJECT INDEX

A CDCC result is accepted only after varying independent truncations.  The
minimum useful record is:
\begin{enumerate}[label=\arabic*.]
  \item projectile Hamiltonian, Pauli prescription, and all fitted structure
  parameters;
  \item Jacobi sets, internal partial waves, Gaussian ranges, and positive-
  energy cutoff;
  \item projectile spins $I$, projectile--target orbital momenta $L$, total
  $J$, and coupling multipoles retained;
  \item matching radius and radial mesh or propagator tolerance;
  \item open and closed channel boundary conditions, including Coulomb
  functions and branch conventions;
  \item optical or folding interactions and the density and energy at which
  they are evaluated;
  \item convergence tables in $E_{\max}$, internal angular momentum, basis
  density, $L_{\max}$, $J_{\max}$, and matching radius;
  \item flux or reaction-cross-section balance;
  \item smoothing angle, scaled-basis size, and independence of the final
  continuous response from these auxiliary choices.
\end{enumerate}
Vary one family of cutoffs at a time and then repeat selected variations
jointly, because Coulomb multipoles can correlate the required excitation
cutoff with the matching radius.  A stable pole, a converged integrated
strength, and a converged differential line shape are three different claims
and require three different tests.

\begin{checkpointbox}
Given a cluster Hamiltonian and fragment--target interactions, the reader
should now be able to construct $P_N$, derive
Eq.~\eqref{eq:cdcc-coupled-radial}, distinguish open from closed channels,
calculate the coupling form factors, and explain why pseudostate strengths
must be smoothed before comparison with an energy-differential measurement.
The next Laboratory performs these steps for $^6$Li and for a
$p+p+{}^7$Be description of $^9$C.
\end{checkpointbox}
